\section{Conjunction judgments: formal benchmarks}
\label{sec:annex-p19-conjunction-judgments}

These proofs concern the explicitly stated event and payoff models.
They do not prove psychological processing claims or reconstruct
unreported participant responses.

\begin{annexitem}{annex:p19.conjunction-benchmark}{thm:p19-conjunction-benchmark}{A common event space and common denominator}
\begin{proof}
The event \(B\cap E\) is the disjoint union of
\(A\cap B\cap E\) and \((B\setminus A)\cap E\).
Finite additivity and nonnegativity give
\[
 P(B\cap E)=P(A\cap B\cap E)+P((B\setminus A)\cap E)
                         \geq P(A\cap B\cap E).
\]
Division by \(P(E)>0\) proves the conditional inequality
with \(B\). Interchanging \(A,B\) proves the other one.
The displayed decomposition also shows that equality
with \(P(B\mid E)\) holds exactly when the residual
event has probability zero. No product factorization
or independence premise was used.

For finitely many indexed cases, the set of indices
satisfying both attributes is a subset of either
single-attribute set. Its cardinality is therefore
no larger. Division by the same strictly positive
number of cases preserves the inequality. Changing
that denominator changes the question and does not
follow from this argument.

Let \(\mathcal P\) be a nonempty family of probability
measures on the same event space. For \(C\subseteq D\),
each \(Q\in\mathcal P\) satisfies \(Q(C)\leq Q(D)\).
Taking suprema preserves the ordering.
For infima, given \(\varepsilon>0\), choose \(Q\)
with \(Q(D)<\inf_{R\in\mathcal P}R(D)+\varepsilon\).
Then
\[
 \inf_{R\in\mathcal P}R(C)\leq Q(C)\leq Q(D)
                    <\inf_{R\in\mathcal P}R(D)+\varepsilon.
\]
Sending \(\varepsilon\) to zero proves monotonicity
of the lower envelope too. The family must be
the same on both sides. All its values lie in
\([0,1]\), so the nonempty infima and suprema exist.
\end{proof}
\end{annexitem}

\begin{annexitem}{annex:p19.conjunction-interpretations}{thm:p19-conjunction-interpretations}{An eight-atom model for the different conditionals}
\begin{proof}
Begin with four cells of weights
\(.20,.05,.10,.65\) in the order stated in the body.
In each cell \(c\), create two atoms \((c,1),(c,0)\).
Give them weights \(w_c r_c\) and \(w_c(1-r_c)\),
where \(r_c=.9\) for the two \(A\) cells and
\(r_c=.1\) for the other two. Let \(E\) be the
set of atoms whose second coordinate is one.
All eight weights are positive and their sum
is \(\sum_cw_c=1\), so this is a probability model.

Direct summation gives
\[
 \begin{gathered}
 P(A)=.25,\quad P(B)=.30,\quad
 P(A\cap B)=.20,\quad P(A^c\cap B)=.10,\\
 P(E)=.25(.9)+.75(.1)=.30,\\
 P(E\cap A\cap B)=.20(.9)=.18,\quad
 P(E\cap B)=.20(.9)+.10(.1)=.19.
 \end{gathered}
\]
Hence \(P(B\mid A)=.20/.25=.80>.30\),
\(P(E\mid A\cap B)=.18/.20=.90>.19/.30=19/30\),
and the conjunction itself is more probable than
the mutually exclusive \(A^c\cap B\).
Conditioning both targets on the same \(E\), however,
gives \(.18/.30=.60\) and \(.19/.30=19/30\),
in the required inclusion order.
This explicitly separates the inverse, conditional,
exclusive and inclusive questions.

For the final marginal claim, choose any
\(q\in[0,1/2]\) and assign the four \(A,B\) cells
weights
\[
                    (q,\ 1/2-q,\ 1/2-q,\ q).
\]
They are nonnegative and sum to one.
Both marginals are \(1/2\), whereas the intersection
has probability \(q\). Varying \(q\) proves that
the two marginals alone cannot select one joint
probability.
\end{proof}
\end{annexitem}

\begin{annexitem}{annex:p19.conjunction-aggregation}{thm:p19-conjunction-aggregation}{What a pooled reversal does and does not imply}
\begin{proof}
Multiplying \(q_i\leq p_i\) by \(w_i\geq0\) and
summing proves the common-weight inequality.
Its contrapositive says that a reversed mean
requires some individual reversal when the same
respondents and weights underlie both means.

For each of the two numerical respondents, use
three atoms belonging respectively to \(A\cap B\),
\(B\setminus A\), and \(B^c\), with probabilities
\((q_i,p_i-q_i,1-p_i)\).
For \((q_1,p_1)=(.8,.9)\) these are
\((.8,.1,.1)\); for \((q_2,p_2)=(.1,.2)\)
they are \((.1,.1,.8)\). Each respondent thus
has a coherent probability assignment on the
same nested event structure.
Using respondent one only for the conjunction
and respondent two only for the constituent
gives the reversed comparison \(.8>.2\).
No respondent in the construction violates
inclusion. The example establishes the logical
limit of separate-group inference; it does not
assert that the source's experimental groups
were sampled in this biased fashion.
\end{proof}
\end{annexitem}

\begin{annexitem}{annex:p19.conjunction-bets}{thm:p19-conjunction-bets}{Substring inclusion, overlapping occurrences and the exact payout probabilities}
\begin{proof}
If \(\texttt{GRGRRR}\) occurs at positions
\(j,\ldots,j+5\), then \(\texttt{RGRRR}\)
occurs at positions \(j+1,\ldots,j+5\).
If \(\texttt{RGRRRG}\) occurs at positions
\(j,\ldots,j+5\), the same shorter word occurs
at positions \(j,\ldots,j+4\).
These positions remain within the twenty-roll
record. This proves both event inclusions.
The payoff inequality follows by multiplying
their indicator inequality by the same
nonnegative prize and subtracting the same cost.

The twenty-letter word consisting of
\(\texttt{RGRRR}\) followed by fifteen additional
\(\texttt R\)'s has only one \(\texttt G\).
It contains the shorter word and neither longer
word, both of which require two \(\texttt G\)'s.
Under the stated independent positive letter
probabilities this record has positive probability.
At a strictly positive common prize, the expected
payoff inequality is therefore strict.

To calculate the occurrence probabilities, fix
a target word \(s\) of length \(m\). Use states
\(0,\ldots,m\). State \(m\) means that the word
has already occurred and is absorbing. Otherwise,
state \(j\) is the length of the longest prefix
of \(s\) that is a suffix of the observed record.
For a letter \(c\), define \(\delta(j,c)\) as
the length of the longest prefix of \(s\) that is a suffix
of \(s_{1:j}c\), capped at \(m\), and set
\(\delta(m,c)=m\).
An empty prefix is always available.

This transition retains all information required
to detect a new occurrence. Indeed, if a prefix
of length \(k\geq1\) is a suffix after appending
\(c\), its first \(k-1\) letters were a suffix
of the earlier record. Their length is at most
the earlier longest matching length \(j\).
They are therefore also a suffix of \(s_{1:j}\).
Testing the suffixes of \(s_{1:j}c\) consequently
tests every possible new prefix match, including
a completed occurrence. Absorption retains an
earlier occurrence regardless of later letters.

Let \(a_t(j)\) be the integer weight in state
\(j\) after \(t\) letters, assigning weight one
to \(\texttt R\) and two to \(\texttt G\).
Start with \(a_0(0)=1\) and all other values zero.
Then
\[
 a_{t+1}(k)=
 \sum_{j=0}^{m}a_t(j)
 \left(\mathbf1_{\{\delta(j,R)=k\}}
                  +2\mathbf1_{\{\delta(j,G)=k\}}\right).
\]
Induction shows that \(a_t(j)/3^t\) is the
state probability: every existing record
extends with the stated two letter weights.
The weights sum to \(3^t\).
The required probability is \(a_{20}(m)/3^{20}\).

For reproducibility, the full transition rows,
listed for successive states beginning at zero,
are:
\[
\begin{array}{c|l|l}
s&(\delta(j,R))_j&(\delta(j,G))_j\\ \hline
\texttt{RGRRR}&(1,1,3,4,5,5)&(0,2,0,2,2,5)\\
\texttt{GRGRRR}&(0,2,0,4,5,6,6)&(1,1,3,1,3,1,6)\\
\texttt{GRRRRR}&(0,2,3,4,5,6,6)&(1,1,1,1,1,1,6)\\
\texttt{RGRRRG}&(1,1,3,4,5,1,6)&(0,2,0,2,2,6,6)
\end{array}
\]
Twenty applications of the recurrence give
the following vectors; their last entries are
the numerators in the body:
\[
\begin{aligned}
a_{20}^{\texttt{RGRRR}}
 &= (1378849860,698113899,672263286,\\
 &\qquad225962702,75951110,435643544),\\
a_{20}^{\texttt{GRGRRR}}
 &= (281559977,1665601626,558352150,482875020,\\
 &\qquad161872020,54263628,282259980),\\
a_{20}^{\texttt{GRRRRR}}
 &= (1,2235870510,747369326,249818094,\\
 &\qquad83505006,27912654,142308810),\\
a_{20}^{\texttt{RGRRRG}}
 &= (1430395948,734502529,703393886,235761666,\\
 &\qquad79021982,26486370,277222020).
\end{aligned}
\]
This counts the union of overlapping occurrence
events exactly. A single product of letter
probabilities would count only one specified
starting position, and multiplication by the
number of positions would double-count records
with multiple occurrences.

For the unequal-reward example, use three
outcomes of probabilities \(.4,.2,.4\), let
the narrow event be the first outcome and
the broad event the first two. Their probabilities
are .4 and .6, but rewards two and one give
expected money .8 and .6. The event inclusion
still holds; the original same-payoff dominance
premise has been changed.
\end{proof}
\end{annexitem}
